\poczatekwykladu{2}
\section{Własności sigma-ciał}
\label{sec:wlasnosci-sigma-cial}

Niech $\A$ będzie sigma-ciałem podzbiorów zbioru $X$.
Wszystkie dopełnienia w tej części rozumiemy względem $X$.

\begin{fakt}
Zbiór $X$ należy do $\A$.
\end{fakt}

\begin{proof}
Mamy $\varnothing\in\A$, a rodzina $\A$ jest zamknięta na dopełnienia.
Zatem $X=\varnothing^c\in\A$.
\end{proof}

\begin{fakt}[Skończone sumy]
Jeżeli $A_1,\ldots,A_m\in\A$, to
\[
  \bigcup_{k=1}^{m}A_k\in\A.
\]
\end{fakt}

\begin{proof}
Połóżmy $A_k=\varnothing$ dla $k>m$. Wtedy wszystkie zbiory tego ciągu
należą do $\A$ i
\[
  \bigcup_{k=1}^{m}A_k=\bigcup_{k=1}^{\infty}A_k\in\A.
\]
\end{proof}

\begin{fakt}[Przeliczalne i skończone przekroje]
Jeżeli $A_n\in\A$ dla każdego $n\geq1$, to
\[
  \bigcap_{n=1}^{\infty}A_n\in\A.
\]
Rodzina $\A$ jest również zamknięta na skończone przekroje.
\end{fakt}

\begin{proof}
Z praw de Morgana otrzymujemy
\begin{equation}
  \bigcap_{n=1}^{\infty}A_n
  =\left(\bigcup_{n=1}^{\infty}A_n^c\right)^c\in\A.
  \label{eq:przekroj-de-morgan}
\end{equation}
Skończone przekroje otrzymujemy tak samo, stosując prawo de Morgana
do skończonych sum.
\end{proof}

\begin{fakt}[Różnice zbiorów]
Jeżeli $A,B\in\A$, to $A\setminus B\in\A$.
\end{fakt}

\begin{proof}
Wystarczy skorzystać z równości $A\setminus B=A\cap B^c$.
\end{proof}

\begin{fakt}[Różnice symetryczne]
Jeżeli $A,B\in\A$, to $A\mathbin{\triangle}B\in\A$, gdzie
\[
  A\mathbin{\triangle}B=(A\setminus B)\cup(B\setminus A).
\]
\end{fakt}

\begin{proof}
Oba zbiory $A\setminus B$ i $B\setminus A$ należą do $\A$, a rodzina
$\A$ jest zamknięta na sumy dwóch zbiorów.
\end{proof}

\begin{uwaga}[Równoważna postać definicji sigma-ciała]
W definicji sigma-ciała wystarczy założyć, że rodzina $\A$ jest niepusta
oraz zamknięta na dopełnienia i przeliczalne sumy.
Warunek $\varnothing\in\A$ wynika wtedy z pozostałych założeń.
\end{uwaga}

\begin{proof}
Wybierzmy $A\in\A$. Wtedy $A^c\in\A$ oraz
$X=A\cup A^c\in\A$ (sumę dwóch zbiorów można zapisać jako sumę
przeliczalną, powtarzając je w ciągu). Stąd $\varnothing=X^c\in\A$.
\end{proof}

\section{Zbiory przeliczalne i ich dopełnienia}
\label{sec:przeliczalne-dopelnienia}

\begin{przyklad}
Niech $X$ będzie dowolnym zbiorem. Rodzina
\begin{equation}
  \A=\{A\subseteq X:
    A\text{ jest przeliczalny lub }A^c\text{ jest przeliczalny}\}
  \label{eq:sigma-cialo-przeliczalne}
\end{equation}
jest sigma-ciałem. Przez zbiór przeliczalny rozumiemy tutaj także zbiór
skończony, czyli zbiór o mocy nie większej niż $\aleph_0$.
\end{przyklad}

\begin{proof}
Sprawdzamy trzy warunki z definicji~\ref{def:sigma-cialo}.
\begin{enumerate}
  \item Zbiór $\varnothing$ jest przeliczalny, więc $\varnothing\in\A$.
  \item Niech $A\in\A$ i $B=A^c$. Jeżeli $A$ jest przeliczalny,
  to $B^c=A$ jest przeliczalny. Jeżeli $A^c$ jest przeliczalny,
  to przeliczalny jest sam zbiór $B$. W obu przypadkach $B=A^c\in\A$.
  \item Niech $A_n\in\A$ dla każdego $n\geq1$. Rozważamy dwa przypadki.

  Jeżeli wszystkie zbiory $A_n$ są przeliczalne, to ich przeliczalna suma
  $\bigcup_{n=1}^{\infty}A_n$ jest przeliczalna, zatem należy do $\A$.

  Jeżeli pewien zbiór $A_{n_0}$ nie jest przeliczalny, to z $A_{n_0}\in\A$
  wynika, że $A_{n_0}^c$ jest przeliczalny. Ponadto
  \begin{equation}
    \left(\bigcup_{n=1}^{\infty}A_n\right)^c
    =\bigcap_{n=1}^{\infty}A_n^c\subseteq A_{n_0}^c.
    \label{eq:dopelnienie-sumy-przeliczalne}
  \end{equation}
  Dopełnienie sumy jest więc przeliczalne. W tym przypadku również
  $\bigcup_{n=1}^{\infty}A_n\in\A$.
\end{enumerate}
\end{proof}

\section{Przestrzeń mierzalna}
\label{sec:przestrzen-mierzalna}

\begin{definicja}[Przestrzeń mierzalna]
Przestrzenią mierzalną nazywamy parę $(X,\mathcal{F})$, gdzie
$\mathcal{F}$ jest sigma-ciałem podzbiorów zbioru $X$.
\end{definicja}

\section{Przekroje i sigma-ciała generowane}
\label{sec:sigma-ciala-generowane}

\begin{twierdzenie}[Przekrój sigma-ciał]
Niech $(\mathcal{F}_t)_{t\in T}$ będzie rodziną sigma-ciał podzbiorów
ustalonego zbioru $X$. Wówczas
\begin{equation}
  \mathcal{G}=\bigcap_{t\in T}\mathcal{F}_t
  \label{eq:przekroj-sigma-cial}
\end{equation}
jest sigma-ciałem. Dla pustej rodziny przyjmujemy, że jej przekrój
jest równy $\mathcal{P}(X)$.
\label{tw:przekroj-sigma-cial}
\end{twierdzenie}

\begin{proof}
Jeżeli $T=\varnothing$, to $\mathcal{G}=\mathcal{P}(X)$ jest sigma-ciałem.
Załóżmy zatem, że $T\neq\varnothing$.
\begin{enumerate}
  \item Dla każdego $t\in T$ mamy $\varnothing\in\mathcal{F}_t$,
  więc $\varnothing\in\mathcal{G}$.
  \item Jeżeli $A_n\in\mathcal{G}$ dla każdego $n\geq1$, to dla każdego
  $t\in T$ wszystkie zbiory $A_n$ należą do $\mathcal{F}_t$.
  Zatem dla każdego $t\in T$ mamy
  \[
    \bigcup_{n=1}^{\infty}A_n\in\mathcal{F}_t,
  \]
  czyli $\bigcup_{n=1}^{\infty}A_n\in\mathcal{G}$.
  \item Jeżeli $A\in\mathcal{G}$, to $A\in\mathcal{F}_t$ dla każdego
  $t\in T$. Stąd $A^c\in\mathcal{F}_t$ dla każdego $t\in T$,
  a więc $A^c\in\mathcal{G}$.
\end{enumerate}
\end{proof}

\begin{wniosek}
Dla dowolnej rodziny $\mathcal{E}\subseteq\mathcal{P}(X)$ istnieje
najmniejsze, w sensie zawierania, sigma-ciało podzbiorów $X$
zawierające $\mathcal{E}$.
\label{wn:najmniejsze-sigma-cialo}
\end{wniosek}

\begin{proof}
Rozważmy rodzinę wszystkich sigma-ciał zawierających $\mathcal{E}$:
\[
  \mathfrak{S}=\{\mathcal{F}\subseteq\mathcal{P}(X):
    \mathcal{F}\text{ jest sigma-ciałem i }\mathcal{E}\subseteq\mathcal{F}\}.
\]
Rodzina $\mathfrak{S}$ jest niepusta, ponieważ należy do niej
$\mathcal{P}(X)$. Z twierdzenia~\ref{tw:przekroj-sigma-cial} wynika, że
\[
  \mathcal{G}=\bigcap_{\mathcal{F}\in\mathfrak{S}}\mathcal{F}
\]
jest sigma-ciałem. Każdy element $\mathcal{E}$ należy do każdego
$\mathcal{F}\in\mathfrak{S}$, zatem $\mathcal{E}\subseteq\mathcal{G}$.
Jeżeli $\mathcal{F}_0$ jest dowolnym sigma-ciałem zawierającym
$\mathcal{E}$, to $\mathcal{F}_0\in\mathfrak{S}$, więc
$\mathcal{G}\subseteq\mathcal{F}_0$.
\end{proof}

\begin{definicja}[Sigma-ciało generowane]
Najmniejsze sigma-ciało zawierające rodzinę
$\mathcal{E}\subseteq\mathcal{P}(X)$ nazywamy sigma-ciałem
generowanym przez $\mathcal{E}$ i oznaczamy przez $\sigma(\mathcal{E})$.
Zatem
\begin{equation}
  \sigma(\mathcal{E})
  =\bigcap\{\mathcal{F}:\mathcal{F}\text{ jest sigma-ciałem w }X,
     \ \mathcal{E}\subseteq\mathcal{F}\}.
  \label{eq:sigma-generowane}
\end{equation}
\end{definicja}

\section{Sigma-ciało zbiorów borelowskich}
\label{sec:zbiory-borelowskie}

\begin{definicja}[Zbiory borelowskie na prostej]
Sigma-ciało generowane przez wszystkie zbiory otwarte w $\R$
nazywamy sigma-ciałem zbiorów borelowskich i oznaczamy przez
$\mathcal{B}(\R)$. Jego elementy nazywamy zbiorami borelowskimi.
\end{definicja}

\begin{fakt}
Sigma-ciało $\mathcal{B}(\R)$ jest generowane przez rodzinę wszystkich
odcinków otwartych:
\begin{equation}
  \mathcal{B}(\R)
  =\sigma\bigl(\{(a,b):a,b\in\R,\ a<b\}\bigr).
  \label{eq:borel-odcinki}
\end{equation}
\end{fakt}

\begin{proof}
Każdy odcinek otwarty jest zbiorem otwartym, zatem sigma-ciało po prawej
stronie \eqref{eq:borel-odcinki} zawiera się w $\mathcal{B}(\R)$.
Z drugiej strony każdy zbiór otwarty $U\subseteq\R$ jest sumą wszystkich
odcinków $(p,q)\subseteq U$ o wymiernych końcach $p<q$.
Takich odcinków jest co najwyżej przeliczalnie wiele. Zatem każdy zbiór
otwarty należy do sigma-ciała generowanego przez odcinki otwarte,
co daje przeciwną inkluzję.
\end{proof}

\begin{uwaga}
Sama rodzina $\mathcal{J}$ wszystkich odcinków nie jest sigma-ciałem.
Na przykład $(0,1)\in\mathcal{J}$, ale jego dopełnienie w $\R$,
\[
  (0,1)^c=(-\infty,0]\cup[1,+\infty),
\]
nie jest odcinkiem, więc nie należy do $\mathcal{J}$.
\end{uwaga}

\begin{samepage}
\section{Pierścienie i sigma-pierścienie zbiorów}
\label{sec:pierscienie}

\begin{definicja}[Pierścień zbiorów]
Pierścieniem zbiorów w $X$ nazywamy rodzinę
$\mathcal{R}\subseteq\mathcal{P}(X)$ spełniającą warunki:
\begin{enumerate}
  \item $\varnothing\in\mathcal{R}$;
  \item jeżeli $A,B\in\mathcal{R}$, to $A\cup B\in\mathcal{R}$;
  \item jeżeli $A,B\in\mathcal{R}$, to $A\setminus B\in\mathcal{R}$.
\end{enumerate}
\end{definicja}

\begin{definicja}[Sigma-pierścień zbiorów]
Sigma-pierścieniem zbiorów w $X$ nazywamy pierścień
$\mathcal{R}\subseteq\mathcal{P}(X)$ zamknięty na przeliczalne sumy.
Oznacza to, że warunek 2) w definicji pierścienia zastępujemy warunkiem:
jeżeli $A_n\in\mathcal{R}$ dla każdego $n\geq1$, to
\[
  \bigcup_{n=1}^{\infty}A_n\in\mathcal{R}.
\]
\end{definicja}
\end{samepage}
